\clearpage
\VerdierPhysicalPageAnchor{0126}%
\begingroup
\small
\setlength{\emergencystretch}{1.5em}%
\setlength{\abovedisplayskip}{4pt plus 2pt minus 2pt}%
\setlength{\belowdisplayskip}{4pt plus 2pt minus 2pt}%

\VerdierSemanticAnchor{ch2-1-3-7}
\noindent\textbf{1.3.7.}\quad
\VerdierSemanticAnchor{ch2-1-3-7-multi-additive-functor}
Soit
$F\colon \prod_{i\in[n]} C_i\longrightarrow C$
un foncteur multi-additif. Supposons que les produits dénombrables soient
représentables dans $C$. Il résulte immédiatement de
(chap. I, 3.4.3 et 3.4.4) que la $\Pi$-extension du foncteur $F$ aux
complexes à homotopie près :

\VerdierSemanticAnchor{ch2-1-3-7-pi-extension}
\[
  K_{\Pi}F\colon \prod_{i\in[n]} K(C_i)\longrightarrow K(C)
\]

\VerdierSemanticAnchor{ch2-1-3-7-multi-exact}
\noindent est un foncteur multi-exact.

\VerdierSemanticAnchor{ch2-1-3-7-koszul-twist}
Ce foncteur est un $\mathbb Z^n$-foncteur tordu par le cocycle de Koszul
conformément à la règle des signes (chap. I, 1.6.6).

\VerdierSemanticAnchor{ch2-1-3-7-sigma-extension}
Résultats analogues pour le foncteur $K_{\Sigma}F$.

\VerdierSemanticAnchor{ch2-1-3-8}
\medskip
\noindent\textbf{1.3.8.}\quad
\VerdierSemanticAnchor{ch2-1-3-8-canonical-opposite-isomorphism}
Soit $C$ une catégorie additive et
$\iota\colon K(C)^\circ\longrightarrow K(C^\circ)$ l'isomorphisme
canonique $\iota=\iota_1$ (chap. I, 2.5.13).

\VerdierSemanticAnchor{ch2-1-3-8-triangulated-opposite}
Munissons la catégorie $K(C)^\circ$ de la structure triangulée décrite
en (1.1.7).

\VerdierSemanticAnchor{ch2-1-3-8-exactness}
Il résulte de (chap. I, 3.4.1) que $\iota$ est un foncteur exact entre
catégories triangulées. Par suite, $\iota$ est un isomorphisme de
catégories triangulées.

\VerdierSemanticAnchor{ch2-1-3-8-identification}
Nous identifierons dorénavant, sauf mention du contraire, les catégories
triangulées $K(C)^\circ$ et $K(C^\circ)$ à l'aide de l'isomorphisme
$\iota$.
\endgroup
